\section{Fields and multipoint positivity}\label{sec:consequences}

The analytic proof was carried out over $\C$. Its conclusion extends by
algebraic incidence conditions, without embedding another field into
$\C$. The implication from complex Nagata to the corresponding statement
over uncountable algebraically closed characteristic-zero fields is also
discussed by Bansal--Majumder~\cite{BM2025}. We give the incidence
argument with the simultaneous quantifiers needed here.

\begin{corollary}\label{cor:fields}
Let $k$ be an uncountable algebraically closed field of characteristic
zero, and let $r\ge10$. There is a countable union of proper Zariski-closed
subsets of the space of distinct ordered $r$-tuples in $\Pp^2_k$, with
nonempty complement, outside which the conclusion of
Theorem~\ref{thm:main} holds with the same degree and multiplicity
conventions.
\end{corollary}
\begin{proof}
Let $A_{d,\mathbf m}$ be the $\mathbb Q$-defined incidence locus
constructed in the proof of Lemma~\ref{lem:universal}: its points are
the tuples admitting a nonzero degree-$d$ form with multiplicities at
least $\mathbf m$. On affine charts it is defined by the full-column-size
minors of the Taylor-coefficient matrix. The same rational minors define
its base change over any characteristic-zero field.

Whenever $d\sqrt r\le\sum_i m_i$, Theorem~\ref{thm:main} shows that
this locus is proper over $\C$. It is therefore proper over $\mathbb Q$
and remains proper over $k$: a nonzero defining minor on a rational affine
chart stays nonzero after field extension. Exclude these loci for all the
countably many violating pairs $(d,\mathbf m)$. Each tuple remaining
satisfies every asserted inequality. The complement is nonempty: choose
$2r$ affine coordinates in $k$ algebraically independent over $\mathbb Q$.
They avoid all these proper rationally defined loci and all coincidences.
Such coordinates exist by uncountability, since the algebraic closure in
$k$ of a field generated over $\mathbb Q$ by finitely many elements is
countable.
\end{proof}

Let $X$ be the blowup of the plane at a tuple covered by
Theorem~\ref{thm:main} or Corollary~\ref{cor:fields}. Write $H$ for the
pullback of a line and $E_1,\ldots,E_r$ for the exceptional curves.
A real divisor class is \emph{nef} if its intersection with every
irreducible curve is nonnegative, and \emph{strictly nef} if every such
intersection is positive. The multipoint Seshadri constant of the line
bundle $\mathcal O_{\Pp^2}(1)$ at the tuple is
\[
 \varepsilon(\mathcal O(1);p_1,\ldots,p_r)
 =\sup\{t\ge0:H-t\textstyle\sum_iE_i\text{ is nef}\}.
\]

\begin{corollary}\label{cor:sesha}
The class $N_r=\sqrt r\,H-\sum_iE_i$ is strictly nef and satisfies
$N_r^2=0$. In particular,
\[
 \varepsilon(\mathcal O(1);p_1,\ldots,p_r)=\frac1{\sqrt r}.
\]
\end{corollary}
\begin{proof}
Every irreducible curve on $X$ is either an exceptional curve or the
strict transform of an irreducible plane curve. The first kind has
intersection $1$ with $N_r$. For the second kind, the intersection is
$d\sqrt r-\sum_i\mult_{p_i}C>0$. Thus $N_r$ is strictly nef.
The intersection identities $H^2=1$, $H\cdot E_i=0$, and
$E_i\cdot E_j=-\delta_{ij}$ give $N_r^2=0$.

The class $H-r^{-1/2}\sum_iE_i$ is nef, so the Seshadri constant is
at least $r^{-1/2}$. For the opposite bound, fix a positive integer $n$ and put
$e_n=\lceil\sqrt r\,(n+1)\rceil$. There are at most
$r\binom{n+1}{2}$ linear conditions on degree-$e_n$ forms for
multiplicity at least $n$ at all the points, whereas
\[
 \binom{e_n+2}{2}>r\binom{n+1}{2}.
\]
Thus a nonzero such form exists. If $H-t\sum_iE_i$ is nef with
$t\ge0$, intersecting its class with the strict transform of this
form's divisor $D$ gives
\[
 0\le e_n-t\sum_i\mult_{p_i}D\le e_n-trn.
\]
Hence
$t\le e_n/(rn)$, which tends to $1/\sqrt r$ as $n\to\infty$.
This proves the upper bound using only the count of multiplicity
conditions.
\end{proof}
